\section{Measurement Instruments}
\label{sec:method}

The procedures below instantiate the evaluation questions of Sections~\ref{sec:intro}
and~\ref{sec:problem}: how nominal budgets map to realized resources; how much of a result
is discovery versus commit-rule selection; and whether failure-mined commit semantics transfer
across providers. The hierarchy is evaluation question $\rightarrow$ control/instrument
$\rightarrow$ evidence. Names such as Frontier, FTA, Pooled-4, and the adapters identify those
controls in tables.

\subsection{Taxonomy}
\label{sec:method-taxonomy}

\begin{itemize}
\item \textbf{Class~A --- single-generator methods (nominal $B{=}6$ each):}
Frontier; the L1-style, S1-inspired, and TALE-EP-style closed-API adapters.
\item \textbf{Class~B --- selectors over the shared multi-generator pool:}
Pooled-4 (plurality); FTA (failure-trace gated deferral). Zero \emph{additional} model calls after
pool generation; pool construction may cost up to $4B$ nominal generator actions.
\item \textbf{Class~C --- offline oracles:}
coverage upper bounds. Never treated as peer winners.
\end{itemize}
Nominal budget $B$ is the logical entitlement of Section~\ref{sec:problem}; realized
$\mathbf{r}$ is tracked separately.

\subsection{Frontier (Class~A diagnostic)}
\label{sec:method-frontier}

Frontier is the Class~A accounting diagnostic, not a production algorithm recommendation: a
reserved-attempt full-expansion policy with fixed parameters across cells. It runs two direct-chain
attempts, forms a gold-free incumbent by plurality, then uses the remaining nominal budget for
challenger expansion/verification before applying a fixed gold-free override rule. On preserved
validated records, reconstructed successful calls average $5.38$ of nominal $6$
(Section~\ref{sec:compute}).

The implementation template retained an early-commit gate with sentinel thresholds
$(2.0,2.0,-1.0)$. Because top support is a vote-share in $[0,1]$, the gate never fires; reported
runs therefore do not evaluate adaptive early stopping. Full pseudocode and action semantics appear
in \ref{app:frontier-pseudo}.

\subsection{Class~B selectors}
\label{sec:method-classb}

\paragraph{Pooled-4}
Plurality over $\{\mathrm{Frontier},\mathrm{L1},\mathrm{S1},\mathrm{TALE}\}$ (agreement $\geq 2$;
ties fall back to Frontier). Same extracted answers as FTA; no additional model calls after pool
generation. Its role is selector attribution, not a same-model SC replacement.

\paragraph{FTA (gated deferral diagnostic)}
FTA retains Frontier's answer or defers using Frontier trace metadata. Gates fire in precedence
order; at most one per instance (\ref{app:gates}):
\begin{itemize}
\item \textbf{Low-Depth Guard (FIX-2).} On a single-weak-branch (or support-zero) signature, replace
Frontier by majority among $\{\mathrm{L1},\mathrm{S1},\mathrm{TALE}\}$.
\item \textbf{External-Consensus Fallback (FIX-4).} Else, if Frontier committed via direct incumbent
agreement and all three externals unanimously disagree, adopt that unanimous answer.
\end{itemize}
FTA is the transfer diagnostic for failure-mined commit semantics, not a learned router. Gate rules
were developed on Cohere$\times$GSM8K failure analysis; that cell is development-adjacent
(Section~\ref{sec:provider-analysis}). Runtime features are gold-free (\ref{app:leakage}).

\subsection{Bounded answer repair}
\label{sec:method-scope}

A deterministic support-based re-ranking is computed for every method but overridden by Frontier's
controller-committed answer when present, while L1/S1/TALE typically score the repair layer. Offline
quantification and winner-stability checks appear in \ref{app:repair}; primary tables retain
the historical scoring convention.
